\subsection{A Safe Experiment with Synthetic Host Events}
\label{sec:host-events}

The event interface turns a host notification into an ABCL message.
You can explore that boundary without a timer, an operating-system signal,
or another thread. The following supplemental exercise uses the interface
described by the user's guide and implemented in \filepath{src/signal.lsp};
there is no corresponding runnable event example in \filepath{examples/}.
It changes no original source and installs no operating-system callbacks.

\paragraph{Receive and inspect an event.}
The file \filepath{tutorial/examples/14-host-events.abclr2} defines an
ordinary recording object. Loading it does not install it as a handler:

\wholelisting[firstline=2,lastline=8]{tutorial/examples/14-host-events.abclr2}

The dotted pattern \code{[:signal . parameters]} binds the argument tail
as a list. Each saved record contains that list, the reply destination,
and the sender. A host notification arrives with both reply and sender
\code{nil}: it is not an application call waiting for \code{!}.
The separate \code{:events} query does reply, allowing us to inspect
what the handler has processed.

Start a fresh interpreter at the repository root and save its current
handler before temporarily replacing it:

\begin{codeblock}
\begin{lstlisting}
(inc "tutorial/examples/14-host-events" :use-source t)
[saved-event-handler := *signal-handler*]
[*signal-handler* := event-inbox]
(r2sys:send-signal :tick 7)
[event-inbox <== :events]           ; => (((:TICK 7) NIL NIL))
\end{lstlisting}
\end{codeblock}

The \code{r2sys:} prefix names the exported host interface in the
ABCL/R2 system package. \code{send-signal} accepts an ordinary Lisp
\code{\&rest} argument list and constructs \code{(:signal :tick 7)}.
It submits the notification; returning from that function is not an
acknowledgment from your handler. This finite interactive exercise lets
the runtime dispatch the event before inspecting the inbox. Do not infer
a general ordering guarantee between independent producers, or use a
notification as proof that unrelated work finished.

Passing one list is different from passing its elements as arguments:

\begin{codeblock}
\begin{lstlisting}
[event-inbox <== :clear]
(r2sys:send-signal '(:tick 7))
[event-inbox <== :events]           ; => ((((:TICK 7)) NIL NIL))
\end{lstlisting}
\end{codeblock}

The parameter list now has one element, itself a list. Similarly,
\code{(r2sys:send-signal)} sends an event with an empty parameter list,
while \code{(r2sys:send-signal nil)} sends one with parameter list
\code{(nil)}. Neither call means ``disable events.''

\paragraph{Explore the single pending-request slot.}
The other interface is a \emph{host Lisp variable},
\code{r2sys:*signal-request*}. It is distinct from a variable stored in
an ABCL environment. An ABCL assignment such as
\code{[r2sys:*signal-request* := value]} would not be the host-side
write this interface requires. The teaching file supplies small host
functions to make that boundary explicit:

\wholelisting[firstline=11,lastline=17]{tutorial/examples/14-host-events.abclr2}

These function bodies run in host Lisp. Their \code{setf} writes the
actual host variable; \code{demo-event-slot} reads it. They are teaching
helpers, not additional built-in ABCL/R2 operations.
In the same fresh session, first check that the slot is empty:

\begin{codeblock}
\begin{lstlisting}
(demo-event-slot)                  ; => NIL
[event-inbox <== :clear]
(demo-post-event '(:tick 7))
[event-inbox <== :events]           ; => ((((:TICK 7)) NIL NIL))
(demo-event-slot)                  ; => NIL: consumed by the runtime
\end{lstlisting}
\end{codeblock}

A non-nil slot value becomes \emph{one} argument following \code{:signal}.
A list stored there is not spliced into multiple arguments.
The original polling function makes this explicit:

\wholelisting[firstline=57,lastline=61]{src/signal.lsp}

It saves the value, clears the slot, and forwards \code{(:signal value)}.
The scheduler calls this function at the top of its dispatch loop in
\filepath{src/process.lsp}. The slot is not a background service that
keeps an otherwise idle interpreter polling, and its check does not
interrupt a host Lisp computation halfway through execution.

Two writes before that polling point do not create two pending events:

\begin{codeblock}
\begin{lstlisting}
[event-inbox <== :clear]
(demo-replace-event :first :second)
[event-inbox <== :events]           ; => (((:SECOND) NIL NIL))
\end{lstlisting}
\end{codeblock}

Both writes happen inside one host function, with no ABCL scheduler
step between them. Only the last value remains to be polled. This
demonstrates overwrite, not queueing or reliable event aggregation.
Writing \code{nil} means no pending request; use
\code{(r2sys:send-signal nil)} if you actually need an event whose
argument is \code{nil}. These helpers assume an isolated fresh session;
do not overwrite a pending request owned by another event producer.

\paragraph{Disable delivery, then restore the saved handler.}
The handler binding controls both interfaces:

\begin{codeblock}
\begin{lstlisting}
[event-inbox <== :clear]
[*signal-handler* := nil]
(r2sys:send-signal :dropped)
(demo-post-event :also-dropped)
(demo-event-slot)                  ; => NIL after polling
[*signal-handler* := event-inbox]
[event-inbox <== :events]           ; => NIL: no replay
[*signal-handler* := saved-event-handler]
(eq *signal-handler* saved-event-handler) ; => T
\end{lstlisting}
\end{codeblock}

Disabling delivery does not buffer notifications for later. A polled
slot value is cleared even when no handler is installed. Conversely,
changing the handler binding is not a cancellation of messages already
sent to a previous handler. Always restore the object you saved, rather
than assuming that the default handler is still the desired one.
The normal default handler prints an ignored-signal diagnostic; ``ignored''
does not necessarily mean silent.

No step in this exercise enables keyboard interrupts or timer delivery.
The historical keyboard-interrupt wrapper in \filepath{src/signal.lsp}
contains host-specific branches; its fallback simply executes its body.
Do not infer an SBCL operating-system hook from the presence of this API.
Nor is this an Express-message experiment: the original Interrupt and
Express examples have their own reflective execution protocols.

\paragraph{Evidence and exercise.}
Run \code{make -C tutorial host-events-check}. Twenty-five probes in two
fresh processes check handler preservation, direct argument shapes,
nil sender/reply fields, slot consumption and overwrite, disabled delivery,
absence of replay, and restoration. Inputs, copied teaching source and
transcripts remain in \filepath{build/host-event-tests/}. Tests install
no operating-system callbacks and modify no runtime or original sources.

Why does a list-valued slot notification differ from
\code{(r2sys:send-signal :tick 7)}? \emph{Answer:} The slot carries one
value; the function accepts several arguments. Why is a zero-valued
slot request still delivered? \emph{Answer:} Only \code{nil} is false
in Lisp; zero is non-nil. Why should the signal handler not rely on
\code{!} to acknowledge the host? \emph{Answer:} The notification has
no application reply destination. Define a separate acknowledgment
protocol if your application needs one.
